\documentclass[11pt,letterpaper]{article}
\usepackage[T1]{fontenc}
\usepackage[utf8]{inputenc}
\usepackage{lmodern,microtype}
\usepackage[margin=1in]{geometry}
\usepackage{amsmath,amssymb,amsthm,mathtools}
\usepackage{booktabs,array,longtable}
\usepackage{xcolor}
\usepackage{enumitem}
\usepackage{fancyhdr}
\usepackage{listings}
\usepackage[unicode,colorlinks=true,linkcolor=black,citecolor=black,urlcolor=blue!45!black]{hyperref}
\usepackage[nameinlink,noabbrev]{cleveref}
\hypersetup{
 pdftitle={A First-Triple Criterion for Log-Concavity of MDS Weight Distributions},
 pdfauthor={Research note prepared for Vladimir Reshetnikov},
 pdfsubject={Dimension-free criterion, boundary cases, and exact verification},
 pdfkeywords={MDS codes, log-concavity, alternating sums, weight distribution}
}
\setlength{\headheight}{14pt}
\pagestyle{fancy}
\fancyhf{}
\fancyhead[L]{\small MDS weight distributions}
\fancyhead[R]{\small A first-triple criterion}
\fancyfoot[C]{\thepage}
\setlength{\parskip}{4pt}
\setlength{\parindent}{1.2em}
\setlist{nosep}
\numberwithin{equation}{section}
\newtheorem{theorem}{Theorem}[section]
\newtheorem{lemma}[theorem]{Lemma}
\newtheorem{proposition}[theorem]{Proposition}
\newtheorem{corollary}[theorem]{Corollary}
\theoremstyle{definition}
\newtheorem{definition}[theorem]{Definition}
\newtheorem{example}[theorem]{Example}
\newtheorem{remark}[theorem]{Remark}
\newcommand{\F}{\mathbb F}
\newcommand{\N}{\mathbb N}
\newcommand{\LC}{\operatorname{LC}}
\newcommand{\wt}{\operatorname{wt}}
\lstset{basicstyle=\small\ttfamily,columns=fullflexible,breaklines=true,
  frame=single,rulecolor=\color{black!20},showstringspaces=false,
  xleftmargin=0.5em,xrightmargin=0.5em}
\title{\textbf{A First-Triple Criterion for Log-Concavity\ of MDS Weight Distributions}\\[0.5em]
\large Removing the dimension restriction by alternating-sum contraction}
\author{Research note prepared for Vladimir Reshetnikov}
\date{20 September 2026}

\begin{document}
\maketitle
\begin{abstract}
Let $C$ be a linear $[n,k,d]_q$ maximum-distance-separable code, with
$d=n-k+1$, and let $A_w$ count its words of Hamming weight $w$.
Log-concavity is understood after zero coefficients have been deleted,
with $A_0=1$ retained. We prove that, whenever $d\ge2$, $k\ge3$, and
$q>d$, this entire distribution is log-concave if and only if
$A_{d+1}^{\,2}\ge A_dA_{d+2}$. There is no restriction on $k/n$ or on
$n$ relative to $q$. The criterion is the single quadratic inequality
\[
 (n+1)q^2-Bq+\frac{dB}{2}\ge0,
 \qquad B=k(d^2+2d-1)+2.
\]
This removes the dimension restriction in the MDS result of Shi, Wang,
An, and Kim and answers the high-dimension extension discussed in
Remark~2 of their 2024 preprint. The proof normalizes the MDS weight
formula into alternating negative-binomial partial sums. Their
log-convexity defect contracts geometrically, fast enough to be
absorbed by the binomial coefficient's log-concavity at every later
weight. We also classify the boundary $q=d$, obtain an exact
classification of single-parity-check codes, and derive sharp
fixed-distance and fixed-rate thresholds. The accompanying standard-library
Python programs reproduce exact checks on 154,888 main-case parameter
triples, additional boundary cases, and 64 explicitly enumerated small
codes. The results are proved for formal weight arrays as well as for
existing codes; none is an assertion that arbitrary MDS parameters are
realizable.
\end{abstract}

\clearpage
\begingroup
\setlength{\parskip}{0pt}
\tableofcontents
\endgroup
\newpage

\section{Problem selection and the precise research target}

\subsection{The random draw}
The area was selected before choosing a problem. A single call to
Python's \texttt{secrets.randbelow(24)} returned the zero-based index
$16$, selecting \emph{coding theory}. The ordered list of candidate
areas, the method, and the result are preserved in
\texttt{data/random\_selection.json}. No area was redrawn.
This records the actual outcome, not a pseudorandom seed reconstructed
after the choice.

\subsection{The question from the literature}
Shi, Wang, An, and Kim establish an MDS log-concavity criterion under
$3\le k\le n/2+3$; their Remark~2 discusses extending it to larger
$k$ \cite[Theorem~9 and Remark~2]{SWAK}. The question addressed here is:
\begin{quote}
Does the same first-triple quadratic criterion suffice without the
restriction $k\le n/2+3$?
\end{quote}
The answer proved below is affirmative for every integer $q>d$.
The boundary $q=d$, where a weight coefficient vanishes, is treated
separately rather than silently divided out.

The published version's indexed abstract retains the stated dimension
restriction \cite{SWAKjournal}. Literature searches on 20 September 2026
found no subsequent removal of it. This is a bounded literature check,
not a guarantee of priority. The full preprint was examined; the
publisher's full-text PDF could not be retrieved. Consequently the
precise target and numbering above are anchored to the accessible
preprint, not to an assumed identical journal layout.

\subsection{What is, and is not, being solved}
The main theorem supplies a full proof of this extension. It does not
solve the MDS \emph{existence} conjecture. Indeed, it does not use any
conjectural upper bound on the length of an MDS code. The proof is an
inequality argument for the universal MDS weight formula and is valid
even for formal parameters for which no code exists.

Three distinctions matter throughout. First, we delete zero
\emph{coefficients}, not the weight-zero codeword. Second, a formal
weight array is not a construction of a code. Third, an exact finite
check corroborates the proof but cannot replace its universal argument.
The discussion of missing weights by Ezerman, Grassl, and Sol\'e
provides useful background for these distinctions \cite{EGS}.

\section{Definitions and the MDS weight formula}

\begin{definition}
A linear $[n,k,d]_q$ code is a $k$-dimensional subspace of $\F_q^n$
whose smallest nonzero Hamming weight is $d$. It is
\emph{maximum-distance-separable}, or MDS, when $d=n-k+1$.
Write
\[
 A_w=\#\{c\in C:\wt(c)=w\},\qquad
 W_C(z)=\sum_{w=0}^n A_wz^w.
\]
The \emph{compressed weight distribution} is the sequence of positive
$A_w$, in increasing order of $w$, including $A_0=1$.
\end{definition}

A positive sequence $(a_0,\ldots,a_m)$ is log-concave if
$a_i^2\ge a_{i-1}a_{i+1}$ for every internal index. Equivalently, its
successive quotients $a_i/a_{i-1}$ are nonincreasing. The definition
used here agrees with the coding-theoretic convention in
\cite{SWAK}; it is not a statement about a sequence in which the
missing weights remain as zeros.

For example, the binary even-weight code of length four has full
array $(1,0,6,0,1)$ and compressed array $(1,6,1)$. These are different
sequences for purposes of testing adjacent inequalities.

\subsection{Derivation by shortening and inclusion--exclusion}
An MDS generator matrix has every set of $k$ columns independent.
To see this, a dependency causing a rank drop in $k$ selected columns
would produce a nonzero codeword zero in those $k$ positions, of
weight at most $n-k=d-1$, a contradiction. Every smaller set of
columns is therefore independent as well.

Consequently, prescribing zeros in a set $T$ of coordinates leaves
$q^{k-|T|}$ words when $|T|\le k$, and only the zero word when
$|T|\ge k$. Fix a set $S$ of $w\ge d$ positions. Inclusion--exclusion
within $S$, after requiring zeros outside $S$, gives the number of
words with support exactly $S$ as
\begin{equation}\label{eq:ie}
 \sum_{j=0}^{w-d}(-1)^j\binom{w}{j}
       \bigl(q^{w-d+1-j}-1\bigr).
\end{equation}
The subtracted constants are legitimate because the complete
alternating sum of the constant $1$ vanishes when $w>0$.
Using Pascal's identity to telescope \eqref{eq:ie} gives
\begin{equation}\label{eq:mds}
 A_0=1,\qquad A_w=0\ (0<w<d),\qquad
 A_w=(q-1)\binom nw
       \sum_{j=0}^{w-d}(-1)^j\binom{w-1}{j}q^{w-d-j}
       \quad(w\ge d).
\end{equation}
This is the usual MDS weight formula, also the starting point of
\cite[Section~5]{SWAK}. Its derivation shows that all linear MDS codes
with the same parameters have the same weight distribution.

\subsection{Formal arrays and a basic feasibility constraint}
Formula \eqref{eq:mds} defines an array for any integers
$1\le k\le n$ and $q\ge2$, whether or not a code exists. We call this
the \emph{formal MDS array}. All statements about it below concern
its numerical coefficients only.

The first three potentially nonzero weights are
\begin{align}
 A_d&=(q-1)\binom nd,\label{eq:first0}\\
 A_{d+1}&=(q-1)\binom n{d+1}(q-d),\label{eq:first1}\\
 A_{d+2}&=(q-1)\binom n{d+2}
     \left(q^2-(d+1)q+\frac{d(d+1)}2\right).
     \label{eq:first2}
\end{align}
If $k\ge2$, \eqref{eq:first1} exists in the array and implies the
necessary feasibility condition $q\ge d$. If $q=d$, that coefficient
is zero. If $q>d$, the proof below shows that \emph{all} coefficients
from $A_d$ through $A_n$ are positive.

\section{Main theorem: one inequality decides everything}

\begin{theorem}[First-triple criterion]\label{thm:main}
Let $n,k,q$ be integers, let $d=n-k+1$, and suppose
\[
 d\ge2,\qquad k\ge3,\qquad q\ge d+1.
\]
For the formal MDS array, the following are equivalent:
\begin{enumerate}[label=(\roman*)]
\item The compressed distribution $(1,A_d,A_{d+1},\ldots,A_n)$ is
log-concave.
\item $A_{d+1}^{\,2}\ge A_dA_{d+2}$.
\item With $B=k(d^2+2d-1)+2$,
\begin{equation}\label{eq:P}
 P_{n,k}(q):=(n+1)q^2-Bq+\frac{dB}{2}\ge0.
\end{equation}
\item $q\ge q_0(n,k)$, where
\begin{equation}\label{eq:root}
 q_0(n,k)=
 \frac{B+\sqrt{B(k-2)(d^2-1)}}{2(n+1)}.
\end{equation}
\end{enumerate}
When these conditions hold, every internal log-concavity inequality
other than the one centered at $A_{d+1}$ is strict.
In particular, the theorem applies to every existing linear MDS code
with these parameters, without a length or dimension restriction.
\end{theorem}

The proof occupies the next two sections. Its only analytical input
is a simple bound on the ratios of successive terms in an alternating
sum. No coding-theoretic classification theorem is used.

\section{Alternating sums and normalized defects}

\subsection{A normalization that fixes the truncation parameter}
Set $s=w-d$ and define
\begin{equation}\label{eq:Bs}
 B_s=\sum_{j=0}^s(-1)^j\binom{d+s-1}{j}q^{s-j}.
\end{equation}
Thus $A_{d+s}=(q-1)\binom n{d+s}B_s$. The symbol $B_s$ denotes a
sequence and is distinct from the coefficient $B$ in \eqref{eq:P}.
Pascal's identity gives
\begin{equation}\label{eq:recurrence}
 B_0=1,\qquad
 B_s=(q-1)B_{s-1}+(-1)^s\binom{d+s-2}{s}.
\end{equation}
For completeness, in \eqref{eq:Bs} expand
$\binom{d+s-1}{j}=\binom{d+s-2}{j}+\binom{d+s-2}{j-1}$.
The first sum is $qB_{s-1}$ plus its last term, and the second is
$-B_{s-1}$.

Divide \eqref{eq:recurrence} by $(q-1)^s$. Define
\begin{equation}\label{eq:terms}
 t_j=\frac{\binom{d+j-2}{j}}{(q-1)^j},\qquad
 T_s=\sum_{j=0}^s(-1)^jt_j.
\end{equation}
Here $t_0=1$, and we obtain the exact identities
\begin{equation}\label{eq:normalized}
 B_s=(q-1)^sT_s,\qquad
 A_{d+s}=(q-1)^{s+1}\binom n{d+s}T_s.
\end{equation}
The crucial feature is that $t_j$ no longer depends on the final
truncation index $s$. It is a single negative-binomial sequence.
The infinite alternating series sums to
$((q-1)/q)^{d-1}$, but its value is not needed for the proof.

\begin{lemma}[Term contraction and positivity]\label{lem:terms}
Under the hypotheses of \cref{thm:main},
\begin{equation}\label{eq:half}
 0<t_1<1,\qquad
 \frac{t_j}{t_{j-1}}=
 \frac{d+j-2}{j(q-1)}\le\frac12\quad(j\ge2).
\end{equation}
All $T_s$ are positive. The odd partial sums are nondecreasing and
satisfy $T_s\ge T_1$ for odd $s\ge1$.
\end{lemma}
\begin{proof}
We have $t_1=(d-1)/(q-1)<1$. Since $q-1\ge d$, the ratio in
\eqref{eq:half} is at most $(d+j-2)/(jd)$. The latter is at most
$1/2$ because $(d-2)(j-2)\ge0$.
Successive pairs satisfy $t_{2r}-t_{2r+1}\ge0$, so
\[
 T_{2r+1}=T_1+\sum_{i=1}^r(t_{2i}-t_{2i+1})\ge T_1>0.
\]
Every even partial sum exceeds the next odd partial sum. This proves
all the assertions and, by \eqref{eq:normalized}, the promised
positivity of all $A_w$ for $w\ge d$.
\end{proof}

\subsection{The competing effects}
For an internal index $1\le s\le k-2$, write $w=d+s$. The binomial
coefficient contributes the factor
\begin{equation}\label{eq:E}
 \frac{\binom nw^2}{\binom n{w-1}\binom n{w+1}}
 =\frac{(w+1)(n-w+1)}{w(n-w)}
 =1+E_s,\qquad
 E_s:=\frac{n+1}{(d+s)(n-d-s)}.
\end{equation}
The powers of $q-1$ in \eqref{eq:normalized} cancel in an adjacent
ratio. Thus
\begin{equation}\label{eq:ratio}
 \frac{A_{d+s}^2}{A_{d+s-1}A_{d+s+1}}
 =(1+E_s)\frac{T_s^2}{T_{s-1}T_{s+1}}.
\end{equation}

For even $s$, $T_s>T_{s-1}$ and $T_s>T_{s+1}$, so the second factor
is greater than one. Every even-offset inequality is automatically
strict.

For odd $s$, instead,
\[
 T_{s-1}=T_s+t_s,\qquad T_{s+1}=T_s+t_{s+1}.
\]
The alternating sum contributes a log-convexity defect
\begin{equation}\label{eq:D}
 D_s:=\frac{T_{s-1}T_{s+1}}{T_s^2}-1
 =\frac{t_s+t_{s+1}}{T_s}
  +\frac{t_st_{s+1}}{T_s^2}>0.
\end{equation}
Consequently
\begin{equation}\label{eq:ratio-odd}
 \frac{A_{d+s}^2}{A_{d+s-1}A_{d+s+1}}
 =\frac{1+E_s}{1+D_s},
 \qquad
 A_{d+s}^2\ge A_{d+s-1}A_{d+s+1}
 \Longleftrightarrow D_s\le E_s.
\end{equation}
It is this comparison of \emph{defect to curvature}, not monotonicity
of the raw adjacent ratios, that works in every dimension.

\begin{lemma}[Uniform defect comparison]\label{lem:defect}
For every odd internal $s\ge1$,
\begin{equation}\label{eq:contraction}
 \frac{D_s}{E_s}\le
 \frac{s}{2^{s-1}}\frac{D_1}{E_1}.
\end{equation}
\end{lemma}
\begin{proof}
Let $c=2^{1-s}$. By \cref{lem:terms},
\[
 t_s\le ct_1,\qquad t_{s+1}\le ct_2,\qquad T_s\ge T_1.
\]
Substituting into \eqref{eq:D} gives
\[
 D_s\le c\frac{t_1+t_2}{T_1}
       +c^2\frac{t_1t_2}{T_1^2}
 \le cD_1.
\]
On the other hand, $d+s\le s(d+1)$ and
$n-d-s\le n-d-1$, with all relevant factors positive. Hence
\[
 \frac{E_s}{E_1}
 =\frac{(d+1)(n-d-1)}{(d+s)(n-d-s)}\ge\frac1s.
\]
Combining these two bounds proves \eqref{eq:contraction}.
\end{proof}

For $s=1$, the multiplier in \eqref{eq:contraction} is one; for every
odd $s\ge3$ it is strictly smaller than one. In particular, the
first odd defect is the only possible obstruction to \emph{all}
log-concavity inequalities holding simultaneously. Later inequalities
can certainly fail when the first fails; the theorem does not say
that a non-log-concave array has only one failure.

\section{Proof of the criterion and the sharp threshold}

\subsection{The boundary involving the zero word}
There is one internal compressed inequality not covered by
\eqref{eq:ratio}: the one centered at $A_d$, whose left neighbor is
$A_0=1$. By \eqref{eq:first0}--\eqref{eq:first1},
\begin{equation}\label{eq:zero-boundary}
 \frac{A_d^2}{A_{d+1}}
 =\frac{q-1}{q-d}\binom nd\frac{d+1}{n-d}>1.
\end{equation}
For example, $\binom nd\ge n$ for $1\le d\le n-1$, and
$(q-1)/(q-d)\ge1$, so the strict inequality is immediate.

\subsection{From the first triple to every triple}
Condition (i) in \cref{thm:main} plainly implies (ii). Conversely,
(ii) is $D_1\le E_1$ by \eqref{eq:ratio-odd}.
For every later odd internal $s$, \cref{lem:defect} gives
\[
 \frac{D_s}{E_s}\le\frac{s}{2^{s-1}}<1.
\]
All even-offset inequalities are strict by the discussion following
\eqref{eq:ratio}, and the zero-word boundary is strict by
\eqref{eq:zero-boundary}. This proves (ii)$\Rightarrow$(i), together
with the theorem's assertion about equality.

\subsection{Reduction to one quadratic}
Substitute \eqref{eq:first0}--\eqref{eq:first2} into (ii) and cancel
positive factors. Since $n-d=k-1$, the result is
\begin{align}
 &(d+2)(k-1)(q-d)^2\notag\\
 &\hspace{2em}-(d+1)(k-2)
       \left(q^2-(d+1)q+\frac{d(d+1)}2\right)\ge0.
       \label{eq:expandedfirst}
\end{align}
The coefficient of $q^2$ is $d+k=n+1$; the coefficient of $q$ is
$-B$; and the constant is $dB/2$. Thus (ii) and (iii) are equivalent.
Notice that $dB$ is even: if $d$ is odd, then $d^2+2d-1$ is even.
The exact test may nevertheless use $2P$ to avoid any division.

A convenient simplification is
\begin{equation}\label{eq:disc}
 B-2d(n+1)=(k-2)(d^2-1).
\end{equation}
The discriminant is therefore $B(k-2)(d^2-1)>0$, and
\begin{equation}\label{eq:atd}
 P_{n,k}(d)=-\frac d2(k-2)(d^2-1)<0.
\end{equation}
The smaller root lies below $d$ and the larger root above it. Since
$q>d$, (iii) is equivalent to being at or above the larger root,
which is exactly \eqref{eq:root}. This proves (iii)$\Leftrightarrow$(iv)
and completes the proof of \cref{thm:main}.

\subsection{Why the dimension restriction disappears}
The earlier argument compares adjacent log-concavity ratios along
odd offsets \cite[proof of Theorem~9]{SWAK}. Their binomial contribution need not increase before
the middle of the weight range. In fact, when $q$ tends to infinity
with $n,d$ fixed, those ratios approach the binomial ratios in
\eqref{eq:E}, which decrease on part of the range. Thus a proof based
on their universal monotonicity would be inappropriate.

The normalized quantity $D_s/E_s$ asks a different question: how
large is the harmful defect relative to the available positive
curvature? The defect loses a factor at least $2^{s-1}$, whereas the
curvature can lose at most a factor $s$ relative to the first
position. The comparison $s\le2^{s-1}$ is independent of $k/n$.
That is the entire mechanism behind the extension.

\section{The boundary and a complete conditional classification}

It is useful to remove the hypothesis $q>d$ carefully. The resulting
classification is conditional on code existence, but the proof again
works for formal arrays. In particular, statements about parameters
with $q=d$ and large $k$ below do not assert that such MDS codes exist.

\begin{theorem}[Boundary cases]\label{thm:boundary}
Let $1\le k\le n$, $d=n-k+1$, and let $q\ge2$ be an integer.
\begin{enumerate}[label=(\roman*)]
\item If $k=1$ or $d=1$, the compressed formal MDS array is
log-concave.
\item For $k\ge2$, $q<d$ is impossible for an MDS code because the
formal coefficient $A_{d+1}$ is negative.
\item If $k=2$ and $q\ge d$, the compressed array is log-concave.
\item If $q=d=2$, the compressed array is log-concave for every $n$.
\item If $q=d\ge3$ and $k\ge3$, the compressed array is log-concave
exactly when
\[
 k\le4,\qquad\text{or}\qquad k=5\ \text{and}\ d\ge9.
\]
\end{enumerate}
Together with \cref{thm:main}, this decides compressed log-concavity
for every existing nonzero linear MDS code.
\end{theorem}

\begin{proof}
If $k=1$, the compressed array is $(1,q-1)$. If $d=1$, the code is
the full space and $A_w=\binom nw(q-1)^w$, a log-concave sequence.
Part (ii) follows from \eqref{eq:first1}. If $k=2$, there are at most
two nonzero positive-weight coefficients. When $q>d$, their only
internal inequality is \eqref{eq:zero-boundary}; when $q=d$, only
$A_d$ remains.

For $d=q=2$, \eqref{eq:recurrence} yields $B_s=1$ for even $s$ and
$B_s=0$ for odd $s$. Thus $A_w=\binom nw$ for even $w$ and $A_w=0$
for odd $w$. The compressed subsequence is log-concave: the quotient
\[
 \frac{\binom n{2r+2}}{\binom n{2r}}
 =\frac{(n-2r)(n-2r-1)}{(2r+2)(2r+1)}
\]
is decreasing in $r$.

Now let $q=d\ge3$. The normalized terms satisfy $T_1=0$, but
$t_j/t_{j-1}<1$ for $j\ge2$. Hence $T_s>0$ for all $s\ge2$:
$T_2=t_2>0$, $T_3=t_2-t_3>0$, and the later alternating pairs
preserve positivity. Therefore the only missing positive weight is
$d+1$. The first support-count factors are
\begin{align}
 B_0&=1,&B_1&=0,&B_2&=\frac{d(d-1)}2,\notag\\
 B_3&=\frac{d(d-1)(d-2)}3,&
 B_4&=\frac{d(d-1)(3d^2-7d+6)}8.
 \label{eq:boundary-B}
\end{align}

We first check the two compressed inequalities affected by deleting
$A_{d+1}$. When $k\ge3$,
\[
 \frac{A_d^2}{A_{d+2}}
 =\frac{2\binom nd^2}{d\binom n{d+2}}>1.
\]
Indeed, $2\le d\le n-2$ gives $\binom nd\ge\binom n2$, and therefore
$\binom nd^2/\binom n{d+2}\ge(d+1)(d+2)/2$.
When $k\ge4$,
\begin{align*}
 \frac{A_{d+2}^2}{A_dA_{d+3}}
 &=\frac{(k-1)(k-2)(d+3)}{(d+1)(d+2)(k-3)}
     \frac{3d(d-1)}{4(d-2)}\\
 &\ge\frac92\frac{d(d-1)(d+3)}{(d-2)(d+1)(d+2)}>1.
\end{align*}
Here $(k-1)(k-2)/(k-3)\ge6$, and the difference between the numerator
and denominator in the last fraction is $d^2+d+4>0$.
Thus $k=3,4$ are settled.

For $k\ge5$, the next internal inequality is centered at $A_{d+3}$.
Using \eqref{eq:boundary-B}, it becomes
\begin{equation}\label{eq:boundaryineq}
 U_d(k-3)-V_d(k-4)\ge0,
\end{equation}
where
\[
 U_d=16(d-2)^2(d+4),\qquad
 V_d=9(3d^2-7d+6)(d+3).
\]
We have
\[
 V_d-U_d=11d^3+18d^2+57d-94>0
\]
for $d\ge3$. The left side of \eqref{eq:boundaryineq} decreases
with $k$. At $k=6$ it equals
\[
 3U_d-2V_d=-6d^3-36d^2-306d+444<0.
\]
Every $k\ge6$ therefore fails this one inequality.

For $k=5$, \eqref{eq:boundaryineq} is the last remaining internal
inequality, and it is
\[
 f(d):=5d^3-18d^2-249d+350\ge0.
\]
The values at $d=3,4,5$ are negative, while
$f(d+1)-f(d)=15d^2-21d-262>0$ for $d\ge5$.
Finally $f(8)=-234$ and $f(9)=296$. Hence, for integer $d\ge3$,
this condition is exactly $d\ge9$. The proof is complete.
\end{proof}

\begin{remark}
The quadratic criterion is intentionally not applied at $q=d$.
For example, the $[6,3,4]_4$ array is
$A_0=1,A_4=45,A_6=18$, which is compressed-log-concave although
$A_5^2<A_4A_6$. Testing the zero coefficient as though it had not
been deleted would answer a different question.
\end{remark}

\section{Exact classification of single-parity-check codes}

For $n\ge2$, let
\[
 C_{n,q}=\{(x_1,\ldots,x_n)\in\F_q^n:x_1+\cdots+x_n=0\}.
\]
This is an $[n,n-1,2]_q$ MDS code for every $n$: its dimension is
$n-1$, it has no word of weight one, and it contains a word of weight
two. Thus there is no length-existence issue in this family.

A direct count gives
\begin{equation}\label{eq:parity}
 A_w=\binom nw\frac{(q-1)^w+(q-1)(-1)^w}{q}.
\end{equation}
For a short derivation, let $N_w$ count $w$-tuples of nonzero elements
summing to zero. A tuple of length $w-1$ with nonzero sum has exactly
one nonzero last coordinate that makes its total zero. Hence
$N_w=(q-1)^{w-1}-N_{w-1}$ with $N_1=0$, which solves to the factor
in \eqref{eq:parity}.

\begin{corollary}\label{cor:parity}
The compressed weight distribution of $C_{n,q}$ is log-concave
exactly in the following cases, where $q$ is a prime power:
\begin{center}
\begin{tabular}{@{}ll@{}}
\toprule
Alphabet size & Permitted lengths $n\ge2$\\
\midrule
$q=2$ & Every length\\
$q=3$ & $n\le3$\\
$q=4$ & $n\le6$\\
$q=5$ & $n\le15$\\
$q\ge7$ & Every length\\
\bottomrule
\end{tabular}
\end{center}
\end{corollary}
\begin{proof}
The binary case is \cref{thm:boundary}. For $q\ge3$ and $n\ge4$,
substitute $d=2$, $k=n-1$, and $B=7n-5$ into \eqref{eq:P}:
\begin{equation}\label{eq:parityP}
 P_{n,n-1}(q)=n(q^2-7q+7)+(q^2+5q-5).
\end{equation}
At $q=3,4,5$, this is respectively $19-5n$, $31-5n$, and
$45-3n$. For $q\ge7$, both coefficients in \eqref{eq:parityP}
are positive. The omitted lengths $n=2,3$ have dimension at most
two and are covered by \cref{thm:boundary}.
\end{proof}

The threshold at $q=5$ exhibits genuine equality. At length $15$,
\[
 A_2=420,\qquad A_3=5460,\qquad A_4=70980,
 \qquad A_3^2=A_2A_4.
\]
Every other internal inequality is strict. At length $16$,
\[
 A_2=480,\qquad A_3=6720,\qquad A_4=94640,
\]
and $A_3^2-A_2A_4=-268800$. The failure occurs near the beginning
of the nonzero weight range, not near the peak of the distribution.

\section{Further consequences and asymptotic regimes}

\subsection{An exact dimension cutoff at fixed distance}
Fix $d\ge2$ and an integer $q>d$, and put
\begin{equation}\label{eq:RS}
 h=d^2+2d-1,\qquad
 R_d(q)=q^2-hq+\frac{dh}{2},\qquad
 S_d(q)=dq^2-2q+d.
\end{equation}
Since $n+1=k+d$, the quadratic in \eqref{eq:P} separates as
\begin{equation}\label{eq:separate}
 P_{n,k}(q)=kR_d(q)+S_d(q).
\end{equation}
Also $S_d(q)>0$ for the stated range.

\begin{corollary}\label{cor:cutoff}
For fixed $d\ge2$ and $q>d$, the formal MDS arrays are log-concave
for every dimension if $R_d(q)\ge0$. If $R_d(q)<0$, the exact
criterion for dimensions $k\ge3$ is
\[
 k\le \left\lfloor\frac{S_d(q)}{-R_d(q)}\right\rfloor.
\]
In particular, at fixed $d,q$, log-concavity cannot fail at one
dimension and then return at a larger dimension.
\end{corollary}
\begin{proof}
Apply \cref{thm:main} to \eqref{eq:separate}. The small dimensions
are already log-concave. In fact $P$ at $k=2$ is
$(d+2)(q-d)^2>0$, so the cutoff is never below two.
\end{proof}

\subsection{A sharp uniform alphabet threshold}
The larger root of $R_d$ is
\begin{equation}\label{eq:qinf}
 Q_\infty(d)=\frac{h+\sqrt{h(d^2-1)}}2.
\end{equation}
As $R_d(d)=-d(d^2-1)/2<0$, the condition $R_d(q)\ge0$ in our
range is exactly $q\ge Q_\infty(d)$.
This is a sharp uniform threshold for the \emph{formal} arrays:
below it, increasing $k$ eventually causes failure; at or above it,
every dimension passes. Existence restrictions may of course prevent
particular failing formal arrays from being realized as codes.

For fixed $d$, the root $q_0(n,k)$ strictly increases with $k$ and
converges to $Q_\infty(d)$. To see strict increase, at the root for
$k$ we have $kR_d+S_d=0$, so $R_d<0$; increasing $k$ makes the
polynomial negative at that old root. The positive root must move
right. Dividing \eqref{eq:separate} by $k$ gives the limit.

\begin{center}
\begin{tabular}{@{}rrr@{}}
\toprule
$d$ & $Q_\infty(d)$, approximately & Smallest prime power at or above it\\
\midrule
% BEGIN GENERATED THRESHOLD ROWS
2 & 5.7913 & 7 \\
3 & 12.2915 & 13 \\
4 & 20.7871 & 23 \\
5 & 31.2829 & 32 \\
6 & 43.7793 & 47 \\
7 & 58.2764 & 59 \\
8 & 74.7739 & 79 \\
9 & 93.2719 & 97 \\
10 & 113.7702 & 121 \\
% END GENERATED THRESHOLD ROWS
\bottomrule
\end{tabular}
\end{center}

Expansion of \eqref{eq:qinf} gives
\begin{equation}\label{eq:qinf-asymp}
 Q_\infty(d)=d^2+\frac32d-\frac54+\frac{1}{4d}
              -\frac{9}{16d^2}+O(d^{-3}).
\end{equation}
This follows by applying the Taylor expansion of $\sqrt{1-x}$ to
$\sqrt{h(h-2d)}$; it is not used to decide any finite case.

\subsection{Fixed positive rate: the threshold is quadratic in length}
\begin{corollary}\label{cor:rate}
Suppose $n\to\infty$ and $k/n\to\rho\in(0,1)$. Then
\begin{equation}\label{eq:rate}
 \frac{q_0(n,k)}{n^2}\longrightarrow\rho(1-\rho)^2.
\end{equation}
In particular, a sequence of such MDS codes over alphabets of size
$q=O(n)$ with $q>d$ eventually fails compressed log-concavity.
\end{corollary}
\begin{proof}
We have $d/n\to1-\rho$ and
$B/n^3\to\rho(1-\rho)^2$. Rewrite \eqref{eq:root}, using
\eqref{eq:disc}, as
\[
 q_0(n,k)=\frac{B}{2(n+1)}
       \left(1+\sqrt{1-\frac{2d(n+1)}B}\right).
\]
The fraction inside the square root tends to zero. Dividing by
$n^2$ proves \eqref{eq:rate}. A linear alphabet size is eventually
smaller than this positive quadratic threshold.
\end{proof}

This is an actual infinite-family consequence, not just a formal
one. Take prime powers $q\to\infty$, set $n=q$, and evaluate all
polynomials of degree less than $k=\lfloor\rho q\rfloor$ at every
element of $\F_q$. A nonzero polynomial has at most $k-1$ roots,
and a product of $k-1$ distinct linear factors attains that number.
The resulting Reed--Solomon code has parameters $[q,k,q-k+1]_q$.
For large $q$, $q>d$, and \cref{cor:rate} proves its failure of
compressed log-concavity.

\subsection{Why a stronger normalized property does not follow}
Normalize the positive-weight coefficients by their binomial factors:
$a_s=A_{d+s}/\binom n{d+s}$. For $d\ge2$, $k\ge3$, $q>d$,
\eqref{eq:normalized} and \eqref{eq:D} give
\[
 \frac{a_1^2}{a_0a_2}=\frac{1}{1+D_1}<1.
\]
Thus the binomial-normalized positive-weight sequence is \emph{never}
log-concave in this range, even when the ordinary weight counts are.
The binomial factor is doing essential work.

Nor does ordinary log-concavity imply real-rootedness of the
compressed generating polynomial. For the $[5,3,3]_7$ parameters,
that polynomial is
\[
 1+60z+120z^2+162z^3.
\]
Its coefficients are log-concave, but its cubic discriminant is
$-74753388$, so it has a nonreal conjugate pair of roots.

\section{Examples, computation, and reproducibility}

\subsection{Selected exact examples}
The following table records values of the integer polynomial $2P$.
For every row, $q>d$ and $k\ge3$, so its sign is the full criterion.
\begin{center}
\begin{tabular}{@{}rrrrrl@{}}
\toprule
$n$ & $k$ & $d$ & $q$ & $2P_{n,k}(q)$ & Compressed log-concavity\\
\midrule
5 & 3 & 3 & 5 & $-8$ & No\\
5 & 3 & 3 & 7 & $104$ & Yes\\
11 & 9 & 3 & 11 & $472$ & Yes\\
4 & 3 & 2 & 3 & $-2$ & No\\
7 & 6 & 2 & 4 & $-8$ & No\\
15 & 14 & 2 & 5 & $0$ & Yes; first-triple equality\\
16 & 15 & 2 & 5 & $-6$ & No\\
\bottomrule
\end{tabular}
\end{center}
The $[11,9,3]_{11}$ example has $9>11/2+3$ and is directly
constructed by polynomial evaluation over $\F_{11}$. Its first
three positive-weight coefficients are
\[
 A_3=1650,\qquad A_4=26400,\qquad A_5=383460,
\]
with $A_4^2-A_3A_5=64251000>0$.

For the boundary example $[6,3,4]_4$, the archive explicitly
enumerates the code generated over $\F_4=\F_2[\omega]$,
$\omega^2=\omega+1$, by
\[
 \begin{pmatrix}
 1&0&0&1&1&1\\
 0&1&0&1&\omega&\omega^2\\
 0&0&1&1&\omega^2&\omega
 \end{pmatrix}.
\]
Its systematic form gives $64$ distinct words. Their exact
weight array is $(1,0,0,0,45,0,18)$.

\subsection{An exact and inexpensive decision procedure}
For the main range, no codeword enumeration and no floating-point
root computation are necessary. Calculate
\[
 d=n-k+1,\qquad B=k(d^2+2d-1)+2,
\]
then test the sign of
\begin{equation}\label{eq:integer-test}
 2(n+1)q^2-2Bq+dB.
\end{equation}
This uses a constant number of arbitrary-precision integer
operations. Its bit cost depends polynomially on the bit lengths of
$n,k,q$; the statement is not that multiplication of unbounded
integers has constant cost. Generating the full array instead takes
$k$ recurrence steps in \eqref{eq:recurrence}, besides computing the
binomial coefficients.

\begin{lstlisting}[language=Python]
def twice_quadratic(n: int, k: int, q: int) -> int:
    d = n - k + 1
    b = k * (d*d + 2*d - 1) + 2
    return 2*(n + 1)*q*q - 2*b*q + d*b
\end{lstlisting}
The complete implementation in \texttt{code/mds.py} validates inputs,
handles all boundary cases, computes the full array, and reports
any directly detected compressed log-concavity failures.

\subsection{Executed verification, with scopes separated}
All reported checks were executed using Python 3.13.5 and exact
integers or \texttt{fractions.Fraction}. No numerical tolerance
was used in a pass/fail decision.

\begin{center}
\begin{tabular}{@{}lr@{}}
\toprule
Check & Number\\
\midrule
Main-case formal parameter triples, $q>d$ & 154,888\\
Formal boundary triples, $q=d\ge3$ & 15,642\\
Binary boundary lengths & 499\\
Degenerate-case parameter checks & 1,633\\
Independent inclusion--exclusion comparisons & 1,815\\
Exact rational checks of \eqref{eq:contraction} & 45,066\\
Explicitly enumerated small codes & 64\\
Codewords enumerated across those codes & 20,798\\
\bottomrule
\end{tabular}
\end{center}

The main sweep uses every $4\le n\le200$ and $2\le d\le n-2$,
with $k=n-d+1$. For each pair it tests the distinct values in
\[
 \{d+1,d+2,2d,3d,n-1,n,n^2,d^2+2d-1\}
\]
that exceed $d$. For every triple it compares the full compressed
array against the theorem, checks positivity on the entire claimed
support, checks that the total is $q^k$, and verifies that any failure
includes the first triple. These triples include non-prime-power
$q$ and unrealizable code parameters deliberately: the theorem is
stronger than a statement restricted to known constructions.

The $q=d\ge3$ sweep uses $3\le d\le200$ and $2\le k\le80$.
The binary checks cover lengths $2$ through $500$.
The rational checks test the central bound directly with exact
fractions for $n\le60$.

The 64 construction checks are separate. They enumerate small
prime-field Reed--Solomon codes, single-parity-check codes over
$\F_2,\F_3,\F_4,\F_5,\F_7$, and the displayed $\F_4$ code.
Their word counts are compared with the universal formula. The
\texttt{verification.json} file records every constructed code's
parameters and the full selected-example arrays.

\subsection{Reproducing the archive}
The source requires an ordinary \LaTeX{} installation with the
standard packages listed in the preamble. From the archive root:
\begin{lstlisting}
latexmk -pdf -interaction=nonstopmode -halt-on-error article.tex
python code/verify.py --max-n 200 --output data/verification.json
python code/mds.py 11 9 11
\end{lstlisting}
Alternatively run \texttt{pdflatex} repeatedly until references stabilize.
The Python programs require version 3.9 or later and only the standard
library. The verification runtime was approximately 21 seconds in
the authoring environment; it naturally varies between machines.
The deterministic digest of the main sweep is included in the JSON
output, so the substantive record can be compared independently of
runtime or Python-version metadata.

\section{Conclusion and status}
The selected extension has an affirmative answer: the first-triple
quadratic criterion does not require a dimension restriction. The
key step is \eqref{eq:contraction}, which compares a geometric
alternating-sum defect with a binomial curvature that can deteriorate
only linearly relative to the first position.

The treatment also closes the zero-coefficient boundary algebraically,
yields an exact classification for every existing linear MDS code,
and gives explicit length thresholds for single-parity-check codes.
At fixed distance there is a finite uniform alphabet threshold; at
fixed positive rate the necessary alphabet threshold grows
quadratically with length. These different regimes explain why
apparently similar families can have very different log-concavity
behavior.

The article supplies a conventional proof and executed exact-arithmetic
checks. It has not been proof-assistant formalized or independently
peer reviewed. No gap is known in the argument presented here, but
novelty and correctness remain subject to independent scrutiny.
The priority claim is limited to the literature check described in
Section~1; the explicit theorem and proof do not depend on that
claim.

\appendix
\section{A compact proof for independent checking}
For a reader auditing the central result without the surrounding
applications, the proof can be compressed to the following chain.
Let $d\ge2$, $k\ge3$, $q\ge d+1$ and $n=d+k-1$. Define
\[
 t_j=\frac{\binom{d+j-2}{j}}{(q-1)^j},\qquad
 T_s=\sum_{j=0}^s(-1)^jt_j.
\]
The MDS recurrence proves
$A_{d+s}=\binom n{d+s}(q-1)^{s+1}T_s$.
Because $t_1<1$ and $t_j/t_{j-1}\le1/2$ for $j\ge2$, all $T_s$
are positive and all odd $T_s\ge T_1$.

Even offsets give strict log-concavity immediately. For odd $s$,
set
\[
 D_s=\frac{t_s+t_{s+1}}{T_s}+\frac{t_st_{s+1}}{T_s^2},\qquad
 E_s=\frac{n+1}{(d+s)(n-d-s)}.
\]
The adjacent inequality is exactly $D_s\le E_s$.
Term contraction gives $D_s\le2^{1-s}D_1$; elementary denominator
comparison gives $E_s\ge E_1/s$. Therefore
\[
 D_1\le E_1\ \Longrightarrow\
 \frac{D_s}{E_s}\le \frac{s}{2^{s-1}}\le1.
\]
The zero-word boundary is strict because
$A_d^2/A_{d+1}>1$. Finally, substituting the first three weights
reduces $D_1\le E_1$ to \eqref{eq:P}; \eqref{eq:atd} identifies the
relevant root without any MDS existence bound.

\section{Artifact inventory and provenance}
\begin{center}
\begin{tabular}{@{}p{0.38\textwidth}p{0.56\textwidth}@{}}
\toprule
File & Purpose\\
\midrule
\texttt{article.tex}, \texttt{article.pdf} & This article and its editable source.\\
\texttt{code/mds.py} & Exact enumerator, full criterion, command-line examples.\\
\texttt{code/verify.py} & Reproducible formal and actual-code checks.\\
\texttt{code/make\_tables.py} & Exact threshold selection and table regeneration.\\
\texttt{data/verification.json} & Executed results, scope, examples, deterministic digest.\\
\texttt{data/random\_selection.json} & The area-selection list, method, and draw.\\
\texttt{data/thresholds.csv} & Numerical values of the exact uniform threshold.\\
\texttt{data/threshold\_rows.tex} & Generated rows mirrored in the article source.\\
\texttt{data/literature\_audit.json} & Sources, exact target, search date, access limitations.\\
\texttt{README.md} & Build and verification instructions.\\

\bottomrule
\end{tabular}
\end{center}
The external source articles are not bundled. Their references are
listed below, and the numerical and mathematical work in this archive
can be inspected without downloading them.

\begin{thebibliography}{9}
\bibitem{SWAK}
Minjia Shi, Xuan Wang, Junmin An, and Jon-Lark Kim.
\emph{Log-Concave Sequences in Coding Theory}.
arXiv:2410.04412v1, 6 October 2024, 31 pages.
Especially Section~5, Theorem~9, and Remark~2 on pages 27--29.
\url{https://arxiv.org/abs/2410.04412}.

\bibitem{SWAKjournal}
Minjia Shi, Xuan Wang, Junmin An, and Jon-Lark Kim.
\emph{Log-Concave Sequences in Coding Theory}.
IEEE Transactions on Information Theory, 71(12):9516--9533, 2025.
Publisher record 11176928. Bibliographic record and indexed abstract
consulted; full journal PDF not retrieved.
\url{https://ieeexplore.ieee.org/document/11176928/}.

\bibitem{EGS}
Martianus Frederic Ezerman, Markus Grassl, and Patrick Sol\'e.
\emph{The Weights in MDS Codes}.
IEEE Transactions on Information Theory, 57(1):392--396, 2011.
Preprint arXiv:0908.1669.
\url{https://arxiv.org/abs/0908.1669}.
\end{thebibliography}
\end{document}
